\subsection{非退化selfの最後の非零corner相}
\label{subsec:cubic-third-root-nonunit-isolated}
\begin{theorem}[比例非退化selfの完成した排除]
\label{thm:cubic-third-root-nonunit-proportional-nondegenerate-empty}
残る $v=1,n=0,T\ne0,\alpha=m=-\omega^2$ も全てH3を与えない。
従って両非零比例の $b^2$ 結合で、selfの二方向が独立となる全領域 $v\ne0$ は排除される。
残る必要候補は $v=0,u=1$、$\alpha\ne0,\alpha^2+\omega\ne0$ と、両 $b^2$ 結合零の領域だけである。
\end{theorem}
\begin{proof}
$T=1$ としてよい。
元middleの $bc,cb$ 射影はrank1である。
その核に入る四つの純 $U^2$ 行を、$c^2$ 行を使って取り出す：
\[
 R=\begin{pmatrix}
 M_{bx}-M_{c^2}\\ M_{xb}-\omega^2M_{c^2}\\M_{yb}\\M_{by}
 \end{pmatrix}_{U^2}.
\]
$\operatorname{rank}M=2$ から $\operatorname{rank}R\le1$ が必要である。
この $4\times4$ 行列の全36個の二次minorを元22tensorから構成し、
$\mathbb Q(\omega)$ 上のGroebner基底11個を計算する。
その線型部分は
\[
 a_1=(1+\omega)a_2,\quad a_5=\omega a_8,\quad a_6=(1+\omega)a_8.
\]
$A=a_0,B_0=a_2,D=a_4,C=a_8,E=a_{15}$、$W=B_0+\omega D-2\omega$ とすると、
残る必要式には
\begin{gather*}
 AW=EW=0,\quad C^2+\omega E(1-D)=0,\\
 B_0^2-\omega B_0=AC,\quad B_0(D-1)+\omega^2AC=0,\\
 D^2-2D+\omega^2B_0-\omega AC=0,\quad AE+B_0C=0
\end{gather*}
がある。公開再現器はこれらを元minor idealのGroebner基底で零へ還元する。

$W\ne0$ なら $A=E=0$、従って $C=0$。
残る三式と $W\ne0$ は $B_0=D=0$ を強制する。
全cube係数が零となり、元middleに純 $c^2$ 関係が現れるので、整域性に反する。
従って $W=0$、$B_0=\omega(2-D)$ である。
必要式は
\[
 AC=\omega^2(2-D)(1-D),\quad
 AE=-\omega(2-D)C,\quad C^2=\omega E(D-1)
\]
となる。

第一の純 $U$ 行を $2\times2$ に並べた行列は
\[
 Q=\begin{pmatrix}-A&\omega^2B_0-1\\1-D&-\omega C\end{pmatrix},
 \qquad \det Q=1-D.
\]
ここで $Q=0$ を捨ててはならない。
$D=1$ では $C=0,AE=0$ であり、$A\ne0$ なら非零rank1の $Q$ により整域性に反する。
$A=0$ なら $Q=0$ で、別の純 $U$ 行が $-x^2+Ey^2$ を与える。
$E=0$ は純 $x^2$ で不可能なので、残余scalarで $E=1$ とできる。
この追加固定点は元22tensorから再構成すると三次商次元22となり、AS4でない。

以下 $D\ne1$ とする。残余scalar作用は $D$ を固定し、$(A,C,E)$ に重み $(-1,1,2)$ で作用する。
これは $p,v,T$ の正規化を保つ対角generator/top変更から直接得られる。
$A\ne0$ では $A=1$ とでき、全係数は
\[
 C=\omega^2(2-D)(1-D),\qquad
 E=-(2-D)^2(1-D),\qquad B_0=\omega(2-D)
\]
という一変数族に入る。
$A=0$ では $D=2,E=\omega^2C^2$ に限り、$C=1$ 又は0の二固定点で覆われる。

一変数族の元関係の像は、全 $D$ について正確に
\begin{gather*}
 xc=\omega^2cx,\quad b^2+yc+\omega^2cy=0,\quad
 \omega^2xb+bx+c^2=0,\\
 \omega^2xb+\omega^2(D-2)yb-bx+(2-D)by=0,\\
 -x^2+(1-D)(xy+yx)-(2-D)(1-D)y^2=0,\\
 xy-yx+\omega^2bc-cb=0
\end{gather*}
となる。
元最高tensorの画像を表示関係で掛け戻し、unit $6\times6$ minorと
係数体以外の分母がないことを確認する。
全次数3,4商から作る正常行列は一般点で $38\times7$、正常消去は7個の係数体定数pivotとなる。
全関係・正常pivotと分母の例外多項式は、平方因子を除いて
\[
 (D-1)(D-2)
\]
だけである。$D=1$ はこの $A=1$ 族では純 $x^2$ 関係で既に除外済み。
$D=2$ を新しく構成すると次数4商次元36となり、AS4値35に反する。
従って一変数族の全値がH3から除外される。

$A=0,D=2,C=1$ の固定点は正常行列rankが全列rankとなり、正常三次は零である。
$C=0$ の固定点は次数4商次元36で非AS4。
先に分離した $Q=0,D=1,E=1$ の点も含めて、全境界の判定が完了した。

公開再現器は \texttt{verify\_cubic\_third\_root\_nonunit\_isolated.py}。
36minor、11個のGroebner基底、必要式の零還元、元画像恒等式、
一変数全pivotと四つのfresh点を標数0で再生する。
前二節の排除と合わせて $v\ne0$ の全領域を除外する。
$v=0$ と両結合零を完成したという意味ではない。
\end{proof}
